\chapter{State of the Art}\label{ch:sota}

Where Chapter~\ref{ch:fundamentals} provides the conceptual \emph{definitional classes}, this state
of the art presents their scientifically substantiated \emph{instances}: for each fundamentals
class, the concrete publications, algorithms, and technologies that realize it. The two chapters
thereby mesh directly (Section~\ref{ssec:measurement-patterns}).

\section{Overview and Dissection Principle}\label{sec:sota-overview}
The analysis rests on 33 in-depth-analysed search-algorithm publications (P01--P33) and a parallel
allocator corpus of 23 works (A01--A23); the axis assignment of each work is persisted as an XML
profile in the \texttt{comdare-cache-engine} (\texttt{algorithm\_profiles/}), and the complete
building-block and allocator matrices are given in Appendix~\ref{app:blocks}. The publications stem
from six thematic \emph{origin clusters} --- the trie family, the hybrid and B\textsuperscript{+}
family, layout theory, prefetching, synchronization, and the hardware-near TU-Dresden works ---;
these clusters serve here only for origin placement, not as a structuring axis.

The core contribution of this thesis is the \emph{dissection} (Section~\ref{ssec:own-framework}):
each overall algorithm (a ``living being'') is decomposed into its orthogonal organs --- the axes
--- and the organ algorithms thus obtained are sorted in \emph{per axis}, with a brief reference to
the origin living being. This axis-wise dissection forms the core of the chapter; the assignment
axis\,$\leftrightarrow$\,algorithm\,$\leftrightarrow$\,publication follows a web-verified provenance
map over roughly 110 organ algorithms.

That the axis scaffold also carries \emph{non-tree-like} methods is demonstrated by a plain hash
table (\texttt{absl::flat\_hash\_map}/SwissTable~\cite{abseil_swisstable}) as a counter-check within
the SearchAlgorithm genus: it satisfies the \emph{point-query subset} of the \texttt{std::map}
interface (\texttt{find}/\texttt{insert}/\texttt{erase}); the order-based operations
(\texttt{lower\_bound}/\texttt{upper\_bound}, ordered range iteration) do not apply, as a hash table
is unordered. Over this subset it populates a \emph{limited axis subset}: the trie-specific axes
(path compression, node type, range index) collapse to pass-through variants, while allocation,
prefetch, concurrency, memory layout, and ISA SIMD (for the SwissTable group lookup) remain shared
--- and thus requires no genus of its own.

\subsection{The Corpus by Origin Clusters}\label{ssec:sota-corpus}
Complementing the axis-wise dissection (Section~\ref{sec:sota-axes}), this overview assigns the
publications to their six origin clusters --- the holistic view of each living being before it is
decomposed into its organs. The \emph{trie family} (cluster~A) comprises ART
(P01)~\cite{leis2013art}, HOT (P02)~\cite{binna2018hot}, Masstree (P03)~\cite{mao2012masstree},
CoCo-Trie (P04)~\cite{boffa2024coco}, START (P05)~\cite{fent2020start}, and the succinct tries LOUDS
(P09)~\cite{jacobson1989louds} and SuRF (P10)~\cite{zhang2018surf}. The
\emph{hybrid and B\textsuperscript{+} family} (cluster~B) unites the B\textsuperscript{2}-tree
(P06)~\cite{schmeisser2022b2tree}, Wormhole (P07)~\cite{wu2019wormhole}, and the cache-sensitive
CSS, CSB\textsuperscript{+}, and Hankins trees
(P11--P13)~\cite{rao1999css,rao2000csb,hankins2003nodesize}. \emph{Layout theory} (cluster~C) ranges
from Samuel (P14)~\cite{samuel2005procaware} and Graefe/Larson (P15)~\cite{graefe2001btree} through
the cache-oblivious designs of Bender et
al.\ (P16/P17)~\cite{bender2000cobtree,bender2005coboblivious} --- theoretically grounded by their
ESA-2002 work on tree layout in the memory hierarchy~\cite{bender2002treelayout}, on cache-oblivious
traversal~\cite{bender2002scanning}, and on the order-maintenance / packed-memory
primitive~\cite{bender2002orderlist} --- to the layout-invariant trees of
Saikkonen/Soisalon-Soininen (P18/P19)~\cite{saikkonen2008multilevel,saikkonen2016layout}.
\emph{Prefetching} (clusters~D/E) comprises \emph{B-Trees Are Back} (Mueller 2025,
P20)~\cite{mueller2025btreesback}, the prefetching B\textsuperscript{+}-trees of Chen et
al.\ (P21/P22)~\cite{chen2001prefetch,chen2002fractal}, Khan (P23)~\cite{khan2010adaptive},
Naderan-Tahan (P24)~\cite{naderan2016adaptivefilter}, Mahling (P25)~\cite{mahling2025hotpath}, as
well as Zhang (P26/P27)~\cite{zhang2024pathprefix,zhang2025hierarchical} and Kuehn
(P28)~\cite{kuehn2023bplustree}. \emph{Synchronization} (cluster~F) carries the optimistic lock
coupling of ART (P08)~\cite{leis2016olc}, RCU (P29)~\cite{mckenney2001rcu}, and hazard pointers
(P30)~\cite{michael2004hazard}; the hardware-near \emph{TU-Dresden line}, finally, Ungethuem
(P31)~\cite{ungethuem2017survey}, Schmidt (P32)~\cite{schmidt2025hbm}, and VAMPIR
(P33)~\cite{berthold2023vampir}. In parallel stands the \emph{allocator corpus} (A01--A23) --- from
the classics dlmalloc~\cite{lea1996dlmalloc}, Hoard~\cite{berger2000hoard}, Slab with its
Magazines/Vmem extension~\cite{bonwick1994slab,bonwick2001vmem}, and Michael's lock-free
allocation~\cite{michael2004lockfree} through the multicore-scalable
mimalloc~\cite{leijen2019mimalloc}, jemalloc~\cite{evans2006jemalloc},
TCMalloc~\cite{ghemawat2005tcmalloc}, snmalloc~\cite{lietar2019snmalloc},
scalloc~\cite{aigner2015scalloc}, rpmalloc~\cite{jansson2017rpmalloc}, and
lrmalloc~\cite{leite2018lrmalloc}, as well as the cache-aware/predictable
CAMA~\cite{herter2011cama}, to NUMAlloc~\cite{yang2023numalloc}, StarMalloc~\cite{starmalloc2024},
and PIM-malloc~\cite{pimmalloc2026}.

\section{Cache Concepts as Instances}\label{sec:sota-cache}
The cache-hierarchy concepts from Section~\ref{sec:cache-basics} appear in the state of the art as
concrete designs and mechanisms; each of the following subsections mirrors a fundamentals definition
from Section~\ref{sec:cache-basics}.

\subsection{Memory and Cache Levels}\label{ssec:sota-levels}
The multi-level hierarchy (Section~\ref{ssec:cache-levels}) is the subject of fundamental
measurements: Graefe and Larson (P15, 2001) study B-tree indexes explicitly against the CPU cache
levels, Bender, Demaine, and Farach-Colton (P16, 2002; P17, 2005) model a multi-level memory model
with a block-transfer parameter and solve the tree-layout problem within it near-optimally even for
an unknown block size~\cite{bender2002treelayout}, and Zhang et al.\ (P27, ASPLOS 2025) study, with
\emph{Hierarchical Prefetching}, a software-/hardware-cooperative \emph{instruction} prefetcher for
server applications; it is relevant as a prefetching instance, not as a data-structure-specific
L1/L2/L3 data prefetch for tree layouts.

\subsection{Cache-Line-Aware Layout}\label{ssec:sota-cacheline}
Cache-line-oriented layout (Section~\ref{ssec:cache-line-sizes}) is the dominant guiding idea of
layout theory: the CSS-tree (P11, Rao/Ross 1999) introduces the cache-line node, the
CSB\textsuperscript{+}-tree (P12, Rao/Ross 2000) the contiguous sibling cluster; Hankins and Patel
(P13, 2003) measure the effect of node size, and Samuel et al.\ (P14, 2005) make the node
processor-aware. The Adaptive Radix Tree (P01~\cite{leis2013art}) graduates the node size
cache-line-orientedly with Node4/16/48/256, and Schmidt et al.\ (P32, 2025) examine strided versus
sequential access on high-bandwidth memory. On the memory-layout axis there are thus, concretely:
dense array-of-structs, cache-line-aligned AoS, column-oriented struct-of-arrays, bit-packed LOUDS
(P09, Jacobson 1989), and the SIMD-tiled AoSoA.

\subsection{Address Translation: TLB and dTLB}\label{ssec:sota-tlb}
The dTLB (Section~\ref{ssec:tlb}) appears in the state of the art mainly as a
\emph{measurement target} --- the measurement architecture maintains a dedicated dTLB-miss category
(Section~\ref{sec:measurement-basics}). On the algorithm side it is addressed by huge pages, which
several allocators (mimalloc, jemalloc, TCMalloc, snmalloc) support, and empirically by Schmidt et
al.\ (P32, 2025) with 2-MiB pages. A \emph{dedicated} dTLB-optimization axis, however, remains
underpopulated in the search-algorithm papers --- an observation to which the research gap
(Section~\ref{sec:gap}) returns.

\subsection{Write Protocols, Update Heuristics, and Inclusion Policy}\label{ssec:sota-update}
The cache-update behaviour introduced in Section~\ref{ssec:coherence} is only partly visible through
the ISA; replacement, inclusion, and concrete coherence are predominantly fixed per
microarchitecture and cache level. \emph{Write protocols:} ARMv8 chooses write-through, write-back,
or non-cacheable, as well as read- or read-write-allocate, through page attributes~\cite{arm_arm},
while x86 uses write-back with write-allocate as the normal case~\cite{intel_sdm}.
\emph{Heuristic replacement/update strategies:} from LRU and pseudo-LRU to
\emph{Re-Reference Interval Prediction} (RRIP/DRRIP), which achieves scan- and thrash-resistant
eviction with two bits per line~\cite{jaleel2010rrip} and builds on the adaptive insertion
policy~\cite{qureshi2007adaptive}. \emph{Inclusion policy:} early Intel client generations
(e.g.\ Nehalem) maintain an inclusive L3, Intel servers since Skylake-SP/Xeon~Scalable a
non-inclusive one, while AMD Zen generations operate the L3 predominantly non-inclusively or
victim-like~\cite{intel_optimization,amd_sog}. Together these three aspects determine which levels a
write updates; since they are microarchitecturally fixed and not selectable at run time, they enter
the system-optimized binary as \emph{measured} platform properties
(Section~\ref{ssec:aware-oblivious}).

\subsection{Cache-Aware and Cache-Oblivious Designs}\label{ssec:sota-aware}
Both philosophies from Section~\ref{ssec:aware-oblivious} are instantiated in the state of the art:
Bender et al.\ (P17, 2005) found the cache-oblivious B-tree on the van Emde Boas layout --- building
on the order-maintenance / packed-memory primitive~\cite{bender2002orderlist} and the
cache-oblivious traversal of ordered sets~\cite{bender2002scanning} ---, P16 (2002) extends the
multi-level layout model, and Saikkonen/Soisalon-Soininen (P18, 2008; P19, 2016) maintain a
cache-sensitive layout invariant even under updates; Graefe/Larson (P15, 2001) compare both
approaches. All practical trie and B\textsuperscript{+} indexes (P01--P14, P20), by contrast, are
designed \emph{cache-aware}. The present thesis positions itself as a third way
(Section~\ref{ssec:aware-oblivious}): measured locality knowledge instead of hard-wired or fully
ignored cache parameters.

\section{Axis Dissection of the Search Methods}\label{sec:sota-axes}
The core of this chapter is the axis-wise \emph{dissection}: each overall algorithm is decomposed
into its organs, and per axis the contributing organ algorithms are listed --- with a brief
reference to the origin living being. The proposed framework maintains
\emph{nineteen composition axes} (T0--T18) as well as three \emph{build axes} (page type, SIMD
extension, hardware profile); many axes are refined into sub-axes. The sub-axis labels
(e.g.\ SA1--SA3 or 6.1--6.5) denote the structuring level of this thesis's building-block matrix
(Appendix~\ref{app:blocks}); in code they are realized as per-topic variant families
(\texttt{axis\_<nn>\_<name>/}). Table~\ref{tab:axes-overview} gives the overview of which
publications contribute an organ per axis; the following subsections elaborate them grouped
thematically. The thematic origin clusters of the prior works (trie family P01--P10,
hybrid/B\textsuperscript{+} family P03/P06/P07/P11--P13, layout theory P14--P19, prefetching
P20--P28, synchronization P08/P29--P30, TU-Dresden hardware P31--P33) thereby appear
\emph{dissolved} into their organ contributions.

\begin{table}[!htbp]
  \caption{Axis dissection of the search methods: nineteen composition axes (T0--T18) and three
  build axes; per axis the sources contributing an organ.}\label{tab:axes-overview}
  \centering\footnotesize
  \begin{tabular}{@{}l l p{6.8cm}@{}}
    \toprule
    \textbf{Axis} & \textbf{Sub-axes} & \textbf{Contributing sources} \\
    \midrule
    T0 search\_algo & SA1--SA3 & P01,P02,P05,P07,P10; k-ary (TUD 2009), interpol., Eytzinger, skip list, B-tree, hash \\
    T1 cache\_traversal & 3.B & cache-line default; P17,P21,P32 \\
    T2 mapping & 3.M & P03,P07; CE baseline \\
    T3 path\_compression & --- & P01,P02,P04,P05,P09 \\
    T4 node\_type & NT1--NT3 & P01,P02,P03,P04,P05,P06,P11,P12,P13 \\
    T5 memory\_layout & HM1--HM4 & P06,P09,P11,P12,P32; default \\
    T6 allocator & 6.1--6.5 & A01--A23; P29,P30 \\
    T7 prefetch & PF1--PF3 & P07,P21,P23,P25,P26,P27 \\
    T8 concurrency & 8.1,8.2 & P08,P29,P30; Dijkstra,Courtois,Herlihy,Michael/Scott \\
    T9 serialization & --- & P01,P09,P10; LZ77 \\
    \midrule
    T10 telemetry & 11.X1--X4 & P01,P02,P16,P26,P28; HdrHistogram (Gil Tene) \\
    T11 value\_handle & --- & P01,P03,P07,P29 \\
    T12 isa & --- & x86-64/AArch64/Power/RISC-V (vendor spec) \\
    T13 index\_organization & --- & Bayer/McCreight,Comer,Oracle \\
    T14 io\_dispatch & IO1--IO3 & Stonebraker,Crotty; baselines \\
    T15 migration\_policy & MG1--MG3 & anti-caching,NVM tiering,LeCaR; none \\
    T16 filter & FT1--FT3 & Bloom,Cuckoo,P10,Xor \\
    T17 queuing\_q1 & --- & Disruptor, Bw-Tree, Driscoll, P29, Pugh, LSM, Lamport, Michael/Scott \\
    T18 queuing\_q2 & --- & eager/lazy/watermark/Kafka/RocksDB \\
    \midrule
    Build page\_type & PG1--PG3 & P01,P02,P04,P05,P06,P10,P11,P12,P13,P20; PRT-ART (26) \\
    Build simd (09b) & --- & SSE2/AVX2/AVX-512/NEON/SVE2/RVV; GPU: CUDA \\
    Build hardware (12) & 12.1--12.5 & HW strategy per paper (Tab.~\ref{tab:hw-sched}) \\
    \bottomrule
  \end{tabular}
\end{table}

\subsection{Search and Navigation Organs}\label{ssec:sota-ax-nav}
The \emph{search axis} (T0, sub-axes SA1--SA3) carries the most organs: as original bindings ART
(P01~\cite{leis2013art}), HOT (P02~\cite{binna2018hot}), START (P05~\cite{fent2020start}), SuRF
(P10~\cite{zhang2018surf}), and Wormhole (P07~\cite{wu2019wormhole}); in addition generic search
\emph{methods}: k-ary search (Schlegel/Gemulla/Lehner~\cite{schlegel2009kary}, TU Dresden),
interpolation search (Perl/Itai/Avni~\cite{perl1978interpolation}), Eytzinger-layout search
(Khuong/Morin~\cite{khuong2017eytzinger}), skip list (Pugh~\cite{pugh1990skiplist}), binary search
tree and B-tree (Bayer/McCreight~\cite{bayer1972btree}), and open-addressing hashing
(Knuth~\cite{knuth1998taocp3}). The \emph{traversal axes} separate the algorithm and cache views: T1
(cache-memory traversal) ranges from cache-line walking through stride prefetch (P21 Chen) and HBM
proximity (P32 Schmidt) to cache-oblivious recursion (P17 Bender~\cite{frigo1999cache}); T2
(mapping) connects the two via linear 1:1 mapping, hash redirect (P07), or permutation index (P03
Masstree) --- an axis sparsely populated in the state of the art, predominantly by a baseline.
\emph{Path compression} (T3) instantiates the byte-wise variant (P01), the single-bit PATRICIA split
(Morrison 1968~\cite{morrison1968patricia}; P02), the collapsed macro node (P04 CoCo), and
multi-byte rewiring (P05). The \emph{node-type axis} (T4, NT1--NT3, 13 building blocks) and the
build axis \emph{page type} (PG1--PG3, 26 building blocks) bundle the adaptive layouts:
Node4/16/48/256 (P01), HOT compound (P02), internal/border (P03), decision/span (P06
B\textsuperscript{2}-tree~\cite{schmeisser2022b2tree}), LOUDS dense/sparse (P10), CSS
node/CSB\textsuperscript{+} node group/Hankins wide (P11/P12/P13), as well as the PRT-ART family.

\subsection{Memory, Allocation, and Value Organs}\label{ssec:sota-ax-mem}
The \emph{memory-layout axis} (T5, HM1--HM4, 8 building blocks) ranges from cache-line-aligned AoS
through pointer-free contiguous arrangement (P11 CSS) and sibling cluster (P12
CSB\textsuperscript{+}) to the embedded tree (P06), bit-packed LOUDS (P09 Jacobson), and the
huge-page variants (P32). The \emph{allocator axis} (T6) is the most densely populated, with the
corpus A01--A23, and is refined into five sub-axes: 6.1 allocation strategy (slab, buddy, region,
pool, object cache), 6.2 reclamation (epoch, RCU P29, hazard pointers P30, QSBR), 6.3 NUMA affinity
(P31 Ungethuem), 6.4 huge-page policy, and 6.5 free-list organization (size-class, best-/first-fit,
segregated). Table~\ref{tab:allocator-profiles} lists the ten runnable allocator profiles with
license and workload affinity. The \emph{value axis} (T11, 5 building blocks) distinguishes inline
value (P01), pointer, redirect-to-node or external pool (P07 Wormhole), versioned pointer (P03
Masstree), and immutable shared-ref (P29 RCU). The \emph{serialization axis} (T9) carries raw
binary, block compression (LZ77, Ziv/Lempel 1977; LZ4/Snappy), succinct LOUDS/FST (P09/P10), and
varint (P01).

\subsection{Hardware-Near Organs}\label{ssec:sota-ax-hw}
The \emph{prefetch axis} (T7, PF1--PF3, 6 building blocks) ranges from no prefetch through
software-fixed (P21 Chen) and adaptively distanced (P23 Khan) to fill-buffer-aware (P25 Mahling),
path-/jump-pointer-oriented (P26 Zhang), and hierarchical instruction-bundle prefetch (P27 Zhang);
the explicit hardware prefetch instruction stems from Wormhole (P07). The \emph{ISA axis} (T12) and
the build axis \emph{SIMD extension} (09b) carry the vendor specifications: x86-64, AArch64, Power,
and RISC-V with SSE2/AVX2/AVX-512, NEON/SVE2, and RVV, as well as --- separately, as a
GPU/accelerator profile (SIMT, not CPU SIMD) --- CUDA. Two overarching strategy axes record what an
algorithm actively \emph{uses} and how it \emph{distributes} work: the \emph{hardware strategy}
(axis 12; sub 12.1 SIMD family, 12.2 cache-level targeting, 12.3 NUMA, 12.4 prefetch hardware, 12.5
atomic family) and the \emph{scheduling strategy} (axis 13; sub 13.1 worker-pool layout, 13.2
SIMD-worker limit --- typically two of $N$ cores, due to hardware ---, 13.3 heterogeneous P-/E-core
dispatch, 13.4 co-routine, 13.5 batch granularity). Table~\ref{tab:hw-sched} assigns both per paper.
The \emph{telemetry axis} (T10, sub 11.X1--X4, 6 building blocks) instantiates the density tracker
(P01), the insert counter (P02), the HdrHistogram (Gil Tene) for p50/p95/p99, the path-read counter
(P26), and the probability hints (P16/P17 Bender). The coherence-friendly counter variants ---
leaf-only, sampling, and offline recompute --- go back to Kuehn (P28, personal communication 2023,
beyond the published DaMoN-2023 state on the cache-/ distribution-aware B\textsuperscript{+}
layout), while the counter in all (including inner) nodes counts as an \emph{anti-pattern}
(cache-line ping-pong).

\begin{table}[!htbp]
  \centering\small
  \caption{Target hardware and scheduling profiles assigned in this thesis per reconstructed SOTA method (selection; the profiles are reconstruction targets of this thesis, not necessarily original-paper properties); the full catalogue is in Appendix~\ref{app:blocks}.}\label{tab:hw-sched}
  \begin{tabular}{@{}l p{4.6cm} p{5.6cm}@{}}
    \toprule
    \textbf{Paper} & \textbf{Hardware strategy (axis 12)} & \textbf{Scheduling strategy (axis 13)} \\
    \midrule
    P01 ART & AVX2, L1-aware, none, CAS & thread-per-core \\
    P02 HOT & AVX2 (bit-vector), L1-aware, CAS & thread-per-core \\
    P03 Masstree & scalar, L1-aware, CAS & thread-per-core + work-stealing \\
    P05 START & AVX2 (multi-byte span), L1-aware, CAS & thread-per-core \\
    P20 B-Trees Are Back & AVX2, HBM-aware, NUMA-aware, CAS & work-stealing + hybrid \\
    P27 hp-soft & PREFETCH (instr.\ bundle), all-cache, CAS & thread-per-core + co-routine \\
    P28 Kuehn & PREFETCH (scalar counter), L1-aware, CAS & thread-per-core + sampling \\
    P31 Ungethuem & AVX-512 (HBM), HBM-aware, NUMA-aware, CAS & hybrid (P-/E-cores) \\
    P32 Schmidt & AVX-512 (SIMD stride), HBM-aware, NUMA-aware, CAS & hybrid \\
    \midrule
    \textbf{PRT-ART} & \textbf{AVX2, L1-aware (TLB-aligned), local NUMA, PREFETCH, CAS (OLC)} & \textbf{thread-per-core + 2-SIMD-worker limit + hybrid + co-routine + micro-batch} \\
    \bottomrule
  \end{tabular}
\end{table}

\subsection{Concurrency, I/O, Migration, and Buffering}\label{ssec:sota-ax-conc}
The \emph{concurrency axis} (T8) divides into 8.1 patterns and 8.2 lock mode. The patterns range
from the serial zero point through coarse locking (Dijkstra~\cite{dijkstra1965concurrent}) and
reader/writer locks (Courtois~\cite{courtois1971readers}) to optimistic lock coupling (P08),
lock-free (Michael/Scott~\cite{michael1996queue}) and wait-free methods
(Herlihy~\cite{herlihy1991waitfree}), as well as RCU (P29) and hazard pointers (P30); the lock mode
distinguishes read-only, read-write, upgradeable, and optimistic validation. The \emph{I/O axis}
(T14, IO1--IO3) instantiates buffered (Stonebraker~\cite{stonebraker1981os}), direct, and
memory-mapped I/O (Crotty~\cite{crotty2022mmap}) alongside the in-memory zero point. The
\emph{migration axis} (T15, MG1--MG3) carries static placement, hot/cold separation (anti-caching
2013), multi-level RAM$\to$NVM$\to$SSD migration (2018), and ML-driven adaptive migration (LeCaR
2018). The \emph{buffering axes} (T17/T18), finally, comprise about a dozen building blocks each:
operation buffers from the FIFO/LIFO queue through the Disruptor ring
buffer~\cite{thompson2011disruptor}, the delta chain (Bw-Tree~\cite{levandoski2013bwtree}), the
copy-on-write buffer (Driscoll~\cite{driscoll1989persistent}), the epoch buffer (P29), the skip list
(Pugh~\cite{pugh1990skiplist}), and the tombstone buffer (LSM, O'Neil~\cite{oneil1996lsm}) to
lock-free SPSC (Lamport~\cite{lamport_spsc}) and MPMC queues
(Michael/Scott~\cite{michael1996queue}); as well as write-back (flush) from eager through lazy,
threshold- and time-window-based (Kafka) to adaptive (RocksDB LSM).

\subsection{Index Organization and Filters}\label{ssec:sota-ax-idx}
The \emph{index-organization axis} (T13) distinguishes the clustered arrangement (storage order
equals index order; Bayer/McCreight~\cite{bayer1972btree}), the unordered heap storage, the
index-organized table (Oracle, VLDB 2000), and the non-clustered secondary index
(Comer~\cite{comer1979btree}). The \emph{filter axis} (T16, FT1--FT3) carries the Bloom
filter~\cite{bloom1970}, the cuckoo filter~\cite{fan2014cuckoo}, the succinct range filter SuRF
(P10), and the Xor filter~\cite{graf2020xor}.

This dissection also reveals the imbalance of the prior works: a few axes are densely and repeatedly
populated (search, node-type, allocator, prefetch axis), others exclusively with baseline/textbook
variants (mapping, cache-memory traversal) or as pure measurement targets (dTLB) --- not a single
method populates all axes jointly. Precisely this incompleteness motivates the research gap
(Section~\ref{sec:gap}).

\section{Workload Frameworks as Instances}\label{sec:sota-workloads}
The terms introduced in Section~\ref{sec:workloads-basics} (Fundamentals) --- load, workload,
workload framework, and the \emph{goal of variation} --- receive their scientific instances here.

\subsection{Benchmark Frameworks}\label{ssec:sota-wl-frameworks}
The loads used in the literature stem from \emph{several}, mutually incomparable benchmark
frameworks; Table~\ref{tab:wl-frameworks} assigns them to the using publications. The most
widespread is the Yahoo!\ Cloud Serving Benchmark (YCSB)~\cite{cooper2010ycsb}; alongside it appear
the index benchmark SOSD~\cite{kipf2019sosd}, the transaction benchmarks of the TPC
family~\cite{tpc_benchmarks}, real system integrations (RocksDB), real string corpora, the CPU suite
SPEC~CPU~\cite{spec_cpu2017}, the scale-out suite CloudSuite~\cite{ferdman2012cloudsuite}, and ---
for the allocator corpus --- mimalloc-bench~\cite{mimalloc_bench} and the classics Larson,
threadtest, shbench, and LADDIS. Because these frameworks produce incomparable loads, the thesis
abstracts them into a common set of load profiles (Section~\ref{ssec:sota-wl-profiles}); the
frameworks not yet realized in the proposed measurement apparatus (TPC, SOSD, SPEC, CloudSuite,
allocator suites) are connected via a dataset loader.

\begin{table}[!htbp]
  \caption{Benchmark frameworks and using publications}\label{tab:wl-frameworks}
  \centering\footnotesize
  \begin{tabular}{@{}p{3.2cm} p{2.8cm} p{6cm}@{}}
    \toprule
    \textbf{Framework} & \textbf{Type} & \textbf{Using publications} \\
    \midrule
    YCSB~\cite{cooper2010ycsb} & KV workload & P02, P03, P07, P20, P26 \\
    SOSD~\cite{kipf2019sosd} & index benchmark & P05 \\
    TPC-C / TPC-DS~\cite{tpc_benchmarks} & OLTP / analytical & P19 \\
    real string corpora & string datasets & P04, P06, P10 \\
    RocksDB integration & real system & P10 (SuRF) \\
    SPEC CPU~\cite{spec_cpu2017} & CPU suite & P23, P24, P27, A18 \\
    CloudSuite~\cite{ferdman2012cloudsuite} & scale-out server & P24, P27 \\
    IBM DB2 & production DBMS & P22 \\
    mimalloc-bench~\cite{mimalloc_bench} & allocator suite & A04, A07, A09, A10, A18 \\
    Larson / threadtest / shbench & allocator classics & A01, A03, P08, P30 \\
    LADDIS & kernel allocator & A02 \\
    AggSum / memory-stream & bandwidth & P32 \\
    custom microbenchmarks & per paper (Zipf/uniform) & P01, P04, P06, P11--P14, P21, P28 \\
    \bottomrule
  \end{tabular}
\end{table}

\subsection{Load Profiles and Workload Affinity of the Methods}\label{ssec:sota-wl-profiles}
From the analysis of the named frameworks, this thesis distills \emph{fourteen} distinct load
profiles (LP01--LP14, Table~\ref{tab:lp-catalog}): they give the heterogeneous loads a common,
uint64-based operation model --- the op-mix names the shares of insert, lookup, delete, scan, and
read-modify-write (RMW) --- over defined key distributions, and each is derived from concrete
publications.

\noindent{\footnotesize\textit{Legend, Tab.~\ref{tab:lp-catalog}, \ref{tab:sota-profiles},
\ref{tab:allocator-profiles}:} Op-mix = shares of insert/lookup/delete/scan/RMW; neg\% = share of
negative (missing-key) lookups; \texttt{expected\_workload} = YCSB workload class (A/B/C/E/F).\par}

\begin{table}[!htbp]
  \caption{Load-profile catalogue (14 profiles, derived from the 33-paper analysis)}\label{tab:lp-catalog}
  \centering\footnotesize
  \begin{tabular}{@{}l p{3.3cm} l l c p{2.4cm}@{}}
    \toprule
    \textbf{LP} & \textbf{Identifier} & \textbf{Op-mix} & \textbf{Distribution} & \textbf{neg\%} & \textbf{Derived from} \\
    \midrule
    LP01 & bulk-insert-dense & 100/0/0/0/0 & uniform-seq & 0 & P01, P02, P03, P05, P12 \\
    LP02 & bulk-insert-sparse & 100/0/0/0/0 & uniform-rand & 0 & P01,P02,P08 \\
    LP03 & insert-value-length-sweep & 100/0/0/0/0 & uniform-seq & 0 & P01,P08 \\
    LP04 & static-read-positive & 0/100/0/0/0 & uniform & 0 & P02, P05, P11, P12, P16, P25 \\
    LP05 & static-read-zipfian & 0/100/0/0/0 & zipfian(.99) & 0 & P10,P20,P25,P28 \\
    LP06 & static-read-negsweep & 0/100/0/0/0 & uniform & 0--100 & P04,P06,P22 \\
    LP07 & static-read-string-corpus & 0/100/0/0/0 & real-corpus & 50 & P04,P06,P10 \\
    LP08 & range-scan-heavy & 0/0/0/100/0 & uniform-seq & 0 & P01,P10,P12,P32 \\
    LP09 & insert-lookup-5050 & 50/50/0/0/0 & uniform & 0 & P03,P29 \\
    LP10 & delete-heavy-post-build & 50/0/50/0/0 & uniform-rand & 0 & P01,P08 \\
    LP11 & mixed-ycsb-oltp & 25/50/0/0/25 & zipfian(.99) & 0 & P10,P19,P20 \\
    LP12 & concurrent-rmw-stress & 25/9/36/9/9 & uniform & 0 & P07,P29 \\
    LP13 & concurrent-read-scaling & 0/100/0/0/0 & uniform & 0 & P07,P08,P25,P28 \\
    LP14 & dynamic-realistic-trace & 9/91/0/0/0 & tpcc-nonuniform & 0 & P19,P20 \\
    \bottomrule
  \end{tabular}
\end{table}

For each SOTA algorithm and each allocator, a profile XML in \texttt{algorithm\_profiles/} records
metadata, the axis assignment (page, node, traversal, value\_handle, concurrency, allocator,
prefetch, telemetry, isa, layout, reclamation), a key-value signature, and an optional
\texttt{<expected\_workload>} tag that overrides the \texttt{ExperimentDriver} heuristic
(Chapter~\ref{ch:methodology}). All 30 SOTA profiles and all 10 allocator profiles are tagged. Three
of the 33 corpus works deliberately have no standalone profile but enter only dissected: the
optimistic lock coupling of ART (P08) as a concurrency organ of the \texttt{art} profile, the
succinct LOUDS representation (P09, Jacobson 1989) as a memory-layout organ (e.g.\ in SuRF, P10),
and the profiling tool VAMPIR (P33) as a measurement technology (Section~\ref{ssec:sota-meas-tech});
the remaining 30 algorithms each carry their own profile XML.

\begin{table}[!htbp]
  \caption{SOTA algorithm profiles with workload affinity}\label{tab:sota-profiles}
  \centering\small
  \begin{tabular}{lll}
    \toprule
    \textbf{Profile} & \textbf{Paper ref.} & \textbf{expected\_workload} \\
    \midrule
    art & P01 (Leis 2013) & YCSB\_C \\
    hot & P02 (Binna 2018) & YCSB\_C \\
    masstree & P03 (Mao 2012) & YCSB\_A \\
    coco\_trie & P04 (Boffa 2024) & YCSB\_C \\
    start & P05 (Fent 2020) & YCSB\_C \\
    b2tree & P06 (Schmeisser 2022) & YCSB\_E \\
    wormhole & P07 (Wu 2019) & YCSB\_A \\
    surf & P10 (Zhang 2018) & YCSB\_A \\
    css\_tree & P11 (Rao/Ross 1999) & YCSB\_C \\
    csb\_tree & P12 (Rao/Ross 2000) & YCSB\_E \\
    hankins & P13 (Hankins 2003) & YCSB\_C \\
    samuel & P14 (Samuel 2005) & YCSB\_C \\
    graefe & P15 (Graefe 2001) & YCSB\_C \\
    bender\_treelayout & P16 (Bender 2002) & YCSB\_C \\
    bender\_cacheoblivious & P17 (Bender 2005) & YCSB\_C \\
    saikkonen\_multilevel & P18 (Saikkonen 2008) & YCSB\_E \\
    saikkonen\_layoutinvariant & P19 (Saikkonen 2016) & YCSB\_E \\
    btreesareback & P20 (Mueller 2025) & YCSB\_E \\
    chen\_prefetch & P21 (Chen 2001) & YCSB\_A \\
    chen\_fractal & P22 (Chen 2002) & YCSB\_A \\
    khan\_adaptive & P23 (Khan 2010) & YCSB\_A \\
    naderan\_tahan & P24 (Naderan-Tahan 2016) & YCSB\_A \\
    mahling & P25 (Mahling 2025) & YCSB\_E \\
    zhang\_fgcs & P26 (Zhang 2024) & YCSB\_A \\
    zhang\_asplos & P27 (Zhang 2025) & YCSB\_A \\
    kuehn & P28 (Kuehn 2023) & YCSB\_C \\
    rcu & P29 (McKenney 2001) & YCSB\_B \\
    hazard\_pointers & P30 (Michael 2004) & YCSB\_B \\
    ungethuem & P31 (Ungethuem 2017) & YCSB\_C \\
    to\_stride & P32 (Schmidt 2025) & YCSB\_C \\
    \bottomrule
  \end{tabular}
\end{table}

The workloads distribute across 13$\times$ YCSB\_C (read-heavy), 9$\times$ YCSB\_A (Zipf-distributed
read+update), 6$\times$ YCSB\_E (range scans), and 2$\times$ YCSB\_B (95/5 read/update). PRT-ART
itself is tagged YCSB\_F (read-modify-write, for density-tracker updates).
Table~\ref{tab:allocator-profiles} lists the ten implemented allocator profiles.

\begin{table}[!htbp]
  \caption{Allocator profiles with license and workload affinity}\label{tab:allocator-profiles}
  \centering\small
  \begin{tabular}{llll}
    \toprule
    \textbf{Profile} & \textbf{Family} & \textbf{License} & \textbf{expected\_workload} \\
    \midrule
    hoard & A01 & Apache-2.0 & YCSB\_A \\
    michael\_lockfree & A03 & BSD-3 & YCSB\_B \\
    mimalloc & A04 & MIT & YCSB\_A \\
    jemalloc & A05 & BSD-2 & YCSB\_B \\
    tcmalloc & A06 & BSD-3 & YCSB\_C \\
    snmalloc & A07 & MIT & YCSB\_A \\
    scalloc & A08 & BSD-3 & YCSB\_A \\
    rpmalloc & A10 & Public Domain & YCSB\_C \\
    lrmalloc & A11 & BSD-3 & YCSB\_B \\
    dlmalloc & A20 & CC0 & YCSB\_C \\
    \bottomrule
  \end{tabular}
\end{table}

\subsection{Bias and Comparability}\label{ssec:sota-wl-bias}
The core finding of the load-profile analysis is a systematic \emph{bias}: each publication chooses
the load profile under which its own design performs best. Static, read-only-oriented tries and
filters (HOT, CoCo, B\textsuperscript{2}-tree, SuRF, CSS) frame ``build-once, 100\,\% positive
lookups'' and avoid updates; compressed tries and filters (CoCo, B\textsuperscript{2}-tree, SuRF)
make the negative-query share the primary axis, because they abort early on misses; Zipf-centred
works (B-Trees Are Back, Mahling, Kuehn) favour cache-aware reordering layouts; write- and
concurrency-heavy works (ART, Masstree, Wormhole, OLC-ART, RCU), conversely, compete only where
static structures fail; and prefetch- and hardware-near works (CSB\textsuperscript{+}, Chen,
Schmidt) measure cache and TLB rather than algorithm logic, via node size and prefetch. Fourteen of
the publications provide no index load profile at all (survey, theory, and hardware works). Only
when \emph{all} load profiles run against \emph{all} living-being binaries (and the container
shells) --- the workload as a dynamic axis (Chapter~\ref{ch:methodology}) --- can the ``winner only
in its self-chosen home profile'' be made visible and bias-reduced comparability be gained. This
incompleteness of the prior works is part of the research gap (Section~\ref{sec:gap}).

\section{Measurement Technologies as Instances}\label{sec:sota-measurement}
The quality criteria and measurement patterns introduced in Section~\ref{sec:measurement-basics}
receive their scientifically substantiated instances here --- the technologies that realize them.

\subsection{Collection Technologies}\label{ssec:sota-meas-tech}
\emph{Hardware performance counters} (PMC, via \texttt{perf}) are nearly universal in the state of
the art and deliver cache-miss rates (L1/L2/L3), dTLB misses, branch misses, and IPC/CPI --- the
majority of the works P01--P28 report such counters. \emph{Profilers} such as OTF2/VAMPIR (P33, TU
Dresden) and VTune collect trace-based runtime profiles, \emph{simulators} such as
gem5~\cite{binkert2011gem5} provide hardware-accurate estimates where real counters are missing (P25
Mahling, P27 Zhang), and the \emph{HdrHistogram} (Gil Tene) collects the latency percentiles
(p50/p95/p99), thereby instantiating the dispersion criterion.

\subsection{Statistical Methodology}\label{ssec:sota-meas-stats}
On the collection the evaluation builds (Section~\ref{ssec:measurement-criteria}, dispersion):
instead of arithmetic means, median and percentiles are reported, differences are tested for
significance with the distribution-free \emph{Mann-Whitney U test} (with Holm correction of the
family-wise error rate), and their strength is quantified via the effect measure
\emph{Cliff's~$\delta$}. This methodology follows the recommendations for scientific benchmarking by
Hoefler and Belli~\cite{hoefler2015benchmarking}.

\subsection{TU-Dresden Hardware Measurement Methodology}\label{ssec:sota-meas-tud}
A methodological line of its own is the hardware-near measurement methodology of TU Dresden (Habich
group): Ungethuem et al.\ (P31) systematize the hardware optimizations for database machines,
Schmidt et al.\ (P32) measure stride versus sequential accesses on HBM with dTLB and L2/L3 counters,
and the VAMPIR tool (P33) profiles non-functional properties of heterogeneous memory. This line
connects the present thesis directly to the database and software-engineering environment of its
university.

\subsection{Measurement Categories and Bridge to the Proposed Apparatus}\label{ssec:sota-meas-bridge}
From these technologies follows the set of measurement categories that the proposed apparatus
maintains: cache-line utilization, cache misses (L1/L2/L3), dTLB misses, memory footprint, branch
misses, IPC/CPI, latency (mean and percentiles p50--p999), throughput, energy, and fill-buffer
occupancy. Of these, the apparatus currently realizes the time- and observer-based categories
(wall-clock, HdrHistogram, per-axis observer, the statistics triad, a two-phase operation loop, and
a conformance gate against \texttt{std::map}); the counter-based PMC categories are foreseen but
their collection is tied to hardware-counter access --- the complete measurement architecture is
developed in Chapter~\ref{ch:methodology}. \emph{Which} of these instances the proposed framework
adopts and which it replaces as obsolete is analysed dialectically (thesis--antithesis--synthesis)
in Chapter~\ref{ch:concept}: each prior work (thesis) meets its limits in the bias-reduced
cross-comparison (antithesis), from which the deeper results of the proposed apparatus emerge as
synthesis. Thus the state of the art connects directly to the research gap (Section~\ref{sec:gap}).

\section{Architectural Design and Software Engineering of Search Structures}\label{sec:sota-design}
Beyond the technical mechanisms, the framework proposed in this thesis is primarily an
\emph{architectural-design} invention. This section places it into the non-technical design and
software-engineering space --- the definitional classes from Section~\ref{sec:cpp23} (design
patterns, metaprogramming) receive their scientific instances here.

\subsection{Data Structures as a Design Space}\label{ssec:sota-design-space}
The idea of regarding data structures not individually but as a \emph{design space} characterizes
Idreos' \emph{Periodic Table of Data Structures}~\cite{idreos2018periodic} and the
\emph{Data Calculator}~\cite{idreos2018datacalculator}: structures arise from the combination of a
few design primitives whose costs can be synthesized. This coincides directly with the axis
dissection of this thesis (Section~\ref{sec:sota-axes}) --- the classical families are mere reading
groupings of orthogonal axis assignments, a concrete method a point in the Cartesian product of the
axes.

\subsection{Self-Designing and Learned Indexes}\label{ssec:sota-design-selftuning}
A second line moves the design into run time: \emph{database cracking}~\cite{idreos2007cracking} and
its robust variant~\cite{halim2012stochastic} build the index adaptively as a by-product of query
processing, while \emph{learned indexes}~\cite{kraska2018learned} replace the structure with a model
of the key distribution; more recent works (VEGA, SIGMOD~2025; AirIndex) continue this in an
actively tuning manner. These approaches show the same trend --- away from the hard-wired single
structure, towards configurable, self-adapting designs --- to which the present thesis gives a
systematic, measurement-driven form with the permutable axis composition.

\subsection{Design-Pattern-Driven Architecture as a Contribution}\label{ssec:sota-design-contribution}
The actual contribution thus lies less in a single structure than in a \emph{tool}: instead of
solving the comparison problem directly, this thesis first constructs the ``shovel'' --- the
permutation and measurement apparatus --- with which arbitrary axis combinations can be generated
and measured fairly. In software terms this apparatus is a composition of named design patterns
(Section~\ref{ssec:classic-patterns}) and two new metaprogramming patterns
(Section~\ref{ssec:new-patterns}). Scientifically it situates itself in the software engineering of
\emph{composition}: invasive software composition from gray-box
components~\cite{assmann2003invasive} (TU Dresden), generative
programming~\cite{czarnecki2000generative}, and feature-oriented software product
lines~\cite{apel2013feature,batory1992hierarchical} --- an axis thereby corresponds to a feature, a
living being to a feature composition.

\emph{Provenance of the axis concept.} The concept of the axis architecture itself is not a new
development of this thesis but adopted from prior work by the author within the brand Comdare, held
by him, respectively the BEP~Venture~UG, where it originally served --- in the context of the
UltiHash database --- the optimization of deduplication procedures. The contribution of the present
thesis is the transfer and shaping of this concept onto cache-aware search structures, as well as
its design-pattern-driven, measurement-oriented realization. \emph{Which} foreign instances the
framework adopts and which it replaces as obsolete is developed dialectically in
Chapter~\ref{ch:concept} (cf.\ Section~\ref{ssec:sota-meas-bridge}).

\section{Research Gap}\label{sec:gap}

The dissection in Section~\ref{sec:sota-axes} exposes the central gap: each of the surveyed works
optimizes only a few axes in isolation, and
\emph{in the corpus studied, not a single method populates all nineteen axes jointly}. Some axes are
populated in the state of the art exclusively with baseline or textbook variants (the mapping
between algorithm and cache traversal), others appear only as a \emph{measurement target} (the
dTLB), and the spatial search class is missing entirely; axes such as I/O, migration, filter, and
buffering are each populated individually, but never together with the search axes in one design.

Added to this is the methodological gap: the results are reported in incomparable settings, and the
the frequent choice of scenario-specific, often favourably tailored load profiles
(Section~\ref{ssec:sota-wl-bias}) hampers a fair cross-comparison. What is missing is a bias-reduced
measurement of \emph{all} load profiles against \emph{all} methods on a common basis.

What is missing --- and what this thesis contributes --- is therefore the systematic, orthogonal
\emph{axis decomposition} of these designs together with a \emph{permuted measurement architecture}
that reconstructs them as explicit configurations and measures them on a common basis. Which of the
instances catalogued here the proposed framework adopts and which it replaces as obsolete is
developed dialectically (thesis--antithesis--synthesis, Section~\ref{ssec:sota-design-contribution})
by the following chapter; from this the deeper results of the proposed apparatus emerge.
